\section{Introduction}
\label{sec:introduction}

Maximum power point tracking (MPPT) is essential to extracting the
available power from a photovoltaic (PV) module or array, whose I-V
characteristic has a single global maximum under uniform illumination but
can develop multiple local maxima under partial shading once bypass diodes
begin conducting. Decades of MPPT research have produced a wide design
space of tracking strategies, but the comparative literature built around
that design space has converged on two recurring, increasingly saturated
patterns: (i) re-comparing perturb-and-observe (P\&O), incremental
conductance (IncCond), and fuzzy logic control against one another, and
(ii) proposing yet another bio-inspired metaheuristic optimizer (grey wolf,
particle swarm, whale, dung beetle, and so on), a pattern reviewers have
grown visibly fatigued by. Neither pattern, on its own, demonstrates
genuine coverage of the MPPT design space, and few published comparisons
release open, reproducible code that a reviewer or reader can independently
re-run.

This paper takes a different approach on both axes. First, rather than
adding a fifth variation on perturbative search or a new metaheuristic, we
compare five algorithms chosen specifically to span four \emph{qualitatively
distinct} MPPT design paradigms:
\begin{itemize}
  \item \textbf{Perturbative} -- P\&O with adaptive step size, and
    IncCond, representing the classical baseline and its refinement;
  \item \textbf{Rule-based} -- a fuzzy logic controller (FLC) with a
    physically-justified 49-rule ($7\times7$) base and centroid
    defuzzification;
  \item \textbf{Learning-based (model-free)} -- tabular Q-learning, trained
    on a subset of irradiance conditions and evaluated separately on
    held-out conditions to directly test generalization rather than
    memorization;
  \item \textbf{Model-based nonlinear control} -- sliding mode control
    (SMC) with a sliding surface on $dP/dV=0$, an exponential reaching law,
    boundary-layer chattering mitigation, and an explicit Lyapunov
    stability argument.
\end{itemize}

Second, every result in this paper is produced by a fully open-source,
pure-Python simulation framework -- no MATLAB/Simulink dependency -- with
pinned dependencies, a continuous-integration test suite, and Monte Carlo
statistical rigor (50 independent runs per algorithm-scenario pair, using
identical random seeds across algorithms for a fair paired comparison). The
underlying photovoltaic model is a two-diode model with bypass-diode
composition, quantitatively validated against a commercial Canadian Solar
CS6P-250P datasheet rather than an idealized or unvalidated model, and the
partial-shading test scenarios are constructed to produce two and three
local maxima, directly probing each algorithm's global-search capability
rather than only its local convergence speed.

\subsection{Contributions}
This paper makes the following contributions:
\begin{enumerate}
  \item A cross-paradigm comparison of five MPPT algorithms spanning
    perturbative, rule-based, learning-based, and model-based control,
    rather than variations on a single paradigm;
  \item An open-source, pure-Python, CI-tested simulation framework
    (two-diode PV model with bypass diodes, boost converter with
    state-space small-signal model, unified algorithm interface, Monte
    Carlo scenario runner) released for full reproducibility;
  \item Explicit, separately-reported train-distribution and held-out
    evaluation for the learning-based algorithm, directly addressing the
    generalization-versus-memorization question that pooled reporting
    would obscure;
  \item A quantitative finding that all five algorithms -- across all four
    paradigms, including the model-based SMC controller with its
    Lyapunov-guaranteed local stability -- fail to reach the true global
    maximum power point in the majority of the complex partial-shading
    patterns tested, with the specific failure mode differing by paradigm
    (Section~\ref{sec:results});
  \item A fuzzy rule-base sensitivity analysis showing a reduced
    $5\times5$ rule base statistically matches the full $7\times7$ base's
    tracking efficiency, a result with direct computational-cost relevance
    given fuzzy logic's substantially higher per-step cost relative to the
    other four algorithms.
\end{enumerate}

The remainder of this paper is organized as follows.
Section~\ref{sec:literature-review} positions this work against prior MPPT
comparison literature, organized by the same four paradigms used above.
Section~\ref{sec:pv-model} and Section~\ref{sec:converter-model} describe
the photovoltaic and boost-converter models and their validation.
Section~\ref{sec:algorithms} details all five algorithm implementations.
Section~\ref{sec:scenarios} defines the test scenarios, Monte Carlo
methodology, and metrics. Section~\ref{sec:results} presents statistical
results and figures. Section~\ref{sec:conclusion} concludes with
limitations and future work.
